\documentclass[10pt]{article}

% Style file switched at submission time; keep neutral here.
\usepackage{amsmath,amssymb,booktabs,graphicx,hyperref,url,microtype}
\usepackage[margin=1in]{geometry}

\title{CodingJEPA: Joint-Embedding Predictive Architecture for Python Code Refactoring}
\author{CodingJEPA Contributors}
\date{2026}

\begin{document}

\maketitle

\begin{abstract}
We introduce CodingJEPA, a Joint-Embedding Predictive Architecture (JEPA) trained to predict
the embedding of a refactored code chunk given the source chunk and a categorical refactoring
intent. Unlike generative approaches, CodingJEPA operates purely in embedding space and
retrieves refactoring candidates from a prebuilt index, avoiding hallucination at the token
level. We train on a corpus of \(\sim\)96,000 human-authored Python refactoring pairs extracted
from 10 open-source repositories (commit cutoff: 2023-12-31). On our CJ-RET-100 benchmark,
CodingJEPA achieves [TODO: R@10 = X.XX], outperforming BM25, MLM-encoder, and frozen
CodeBERT baselines by [TODO: Y.Y\%].
\end{abstract}

\section{Introduction}
\label{sec:intro}

Code refactoring --- transforming source code to improve its internal structure without changing
external behaviour --- is a routine engineering task. Recent large language models can generate
refactored code directly, but they suffer from hallucination: plausible-looking outputs that
subtly change semantics or introduce bugs.

We take a retrieval-based alternative. CodingJEPA predicts \emph{where} a refactored version
lives in a prebuilt embedding index rather than \emph{what} it looks like at the token level.
The prediction target is the L2-normalised embedding of the post-refactoring chunk; the
predictor is an autoregressive Transformer conditioned on the source embedding and a categorical
intent label.

\paragraph{Contributions.}
\begin{itemize}
  \item A JEPA architecture adapted for code: encoder, projector, autoregressive predictor,
        intent embedder, and a sliced isotropic Gaussian regulariser (SigReg).
  \item A reproducible dataset of [TODO: N] Python refactoring pairs across 8 intents, with
        provenance, manifest hashing, and cross-split leakage detection.
  \item A multi-benchmark evaluation harness covering retrieval, execution preservation,
        robustness probes, OOD generalisation, linear probes, and human review.
\end{itemize}

\section{Related Work}
\label{sec:related}

\paragraph{Code representation learning.}
CodeBERT \cite{feng2020codebert} and GraphCodeBERT \cite{guo2021graphcodebert} use masked
language modelling over code–docstring pairs. UniXcoder \cite{guo2022unixcoder} extends this
to a unified architecture. Our encoder shares the masked-token inspiration but is trained
with a predictive objective rather than reconstructive one.

\paragraph{Joint-embedding predictive architectures.}
I-JEPA \cite{assran2023ijepa} introduced the JEPA principle for vision: predict representations
of target regions from context regions, without requiring a decoder. V-JEPA \cite{bardes2023vjepa}
extends this to video. CodingJEPA adapts the JEPA objective to code with categorical
conditioning on refactoring intent.

\paragraph{Code generation and refactoring.}
[TODO: cite relevant LLM/refactoring papers.]

\section{Method}
\label{sec:method}

\subsection{Architecture}

CodingJEPA consists of five modules (Figure~\ref{fig:arch}):
\begin{enumerate}
  \item \textbf{Encoder} $f_\theta$: a 6-layer pre-norm Transformer with RoPE positional
        encoding, 8 attention heads, hidden size 512, GELU activations, and 0.1 dropout.
        Applied to both source and target chunk token sequences.
  \item \textbf{Projector} $g_\theta$: Linear(512, 2048) $\to$ BatchNorm1d $\to$ ReLU $\to$
        Linear(2048, 512).
  \item \textbf{Autoregressive Predictor} $p_\theta$: a 4-layer \texttt{norm\_first}
        TransformerEncoder that predicts the target projection from the source projection and
        intent embedding.
  \item \textbf{Intent Embedder}: $\texttt{nn.Embedding}(9, 512)$; index 8 = \texttt{[I\_NONE]}.
  \item \textbf{SigReg}: sliced isotropic Gaussian regulariser with $K=256$ random projections,
        penalising collapse of the representation space.
\end{enumerate}

\subsection{Training objective}

The primary loss is the cosine distance between the predictor output and the L2-normalised
target projection, summed with the SigReg penalty:
\[
  \mathcal{L} = \bigl(1 - \cos(p_\theta(z_s, e_I),\, g_\theta(f_\theta(x_t)))\bigr)
              + \lambda \cdot \mathcal{L}_{\text{SigReg}}.
\]

\subsection{Training procedure}

Stage A (pretraining): AdamW, lr = 3e-4, weight decay = 0.05, cosine schedule with 5,000
warmup steps, 200,000 total steps, batch 64, gradient accumulation 16.

Stage B (intent fine-tune): 50,000 additional steps on intent-balanced samples.

\section{Data}
\label{sec:data}

\subsection{Corpus construction}

[TODO: fill in after pipeline run completes.]

\subsection{Tokeniser}

SentencePiece BPE with 32,000 vocabulary tokens and 15 special tokens (see
\texttt{codingjepa.data.tokenizer}). Maximum sequence length: 512 BPE tokens.

\section{Experiments}
\label{sec:experiments}

\subsection{Benchmarks}

We evaluate on 12 benchmarks grouped into four families (Table~\ref{tab:main}):

\begin{table}[h]
\centering
\caption{Main results.}
\label{tab:main}
\begin{tabular}{lllll}
\toprule
Benchmark & CodingJEPA & BM25 & MLM-enc & CodeBERT \\
\midrule
CJ-RET-100 R@10 & [TODO] & [TODO] & [TODO] & [TODO] \\
CJ-RET-1k R@10  & [TODO] & [TODO] & [TODO] & [TODO] \\
CJ-EXEC pass rate & [TODO] & --- & [TODO] & [TODO] \\
\bottomrule
\end{tabular}
\end{table}

\subsection{Per-intent breakdown}

[TODO: 8 rows × Retrieval@10.]

\subsection{Robustness probes}

[TODO: CJ-ROB-FMT / CJ-ROB-RENAME / CJ-ROB-DOC rank-change \% and cosine drift.]

\section{Analysis}
\label{sec:analysis}

\subsection{Ablations}

[TODO: matrix from RFC-0005 §D6.]

\subsection{Failure modes}

[TODO: worst-50 retrievals per intent; pointers to \texttt{results/confusions/}.]

\section{Limitations}
\label{sec:limitations}

\begin{itemize}
  \item Corpus contamination: the public pretrained baselines (CodeBERT) are trained on
        data that overlaps with our source repos; we document this in the memo's limitations
        section but cannot fully control for it.
  \item OOD scope: CJ-OOD covers cpython 3.11 $\to$ 3.13 refactors only; cross-language and
        cross-domain generalization is not evaluated.
  \item Commit cutoff: the 2023-12-31 cutoff excludes patterns that emerged in 2024--2025.
  \item Human review (CJ-HUMAN): Cohen's $\kappa$ is not yet available pending annotation.
\end{itemize}

\section{Conclusion}
\label{sec:conclusion}

We have presented CodingJEPA, a retrieval-based code refactoring system trained with a
joint-embedding predictive objective. The architecture avoids token-level generation in favour
of predicting in embedding space and retrieving from a content-addressed index. [TODO: add
conclusion once numbers are available.]

\paragraph{Reproducibility statement.}
Training code, data pipeline, evaluation harness, and model checkpoint are released at
\url{https://github.com/AbdelStark/CodingJEPA}. The corpus is published on HF Hub at
\texttt{CodingJEPA/human-python-refactors}. Running \texttt{make eval} from the repository
root reproduces all numbers in this paper.

\appendix

\section{Hyperparameters}
\label{app:hparams}

[TODO: full hyperparameter table.]

\section{Data statistics}
\label{app:data}

[TODO: per-repo pair counts, intent distribution.]

\section{Extended results}
\label{app:results}

[TODO: full metric tables beyond Table~\ref{tab:main}.]

\section{Human review protocol}
\label{app:human}

[TODO: rating rubric, annotator instructions, inter-annotator agreement.]

\section{Ethical considerations}
\label{app:ethics}

All training data is sourced from permissively-licensed public repositories
(MIT / Apache-2.0 / PSF). The commit cutoff (2023-12-31) limits training to
human-authored commits from the pre-LLM-assisted-commit era (RFC-0002 §D11).
No personal data is collected; email addresses appearing in commits are scrubbed
by the normalization step.

\bibliographystyle{plainnat}
\bibliography{refs}

\end{document}
